%% Vorlage Bachelorarbeit

%% Versionshistorie:

%% v1.0: Erstellung durch Johannes Woske, IT2010, j.woske+latex@gmail.com
%% v2.0: Überarbeitung und Ergänzung durch Anne Traulsen, IT2015, a.traulsen+latex@gmail.com

\documentclass[
	12pt, %Schriftgröße
	a4paper,
	listof=totoc, %Inhaltsverzeichniseinträge für Listen (z.B. Abbildungen)
	bibliography=totoc, %Inhaltsverzeichniseinträge f+r Quellen
	numbers=noenddot, %Entfernt Punkt hinter Gliederungsnummern
	ngerman, %Sprachpaket
	headsepline, %Headertrennlinie
	%footsepline, %Footertrennlinie
	oneside %einseitiges Druckformat %%% Unterdrücken der leeren Seite nach Titelblatt
	]{scrbook} %Dokumentenklasse (Koma-Script)
\usepackage[section]{placeins}
\usepackage[T1]{fontenc}
\usepackage{float}
\usepackage[utf8]{inputenc}
\usepackage[ngerman]{babel}
\usepackage{url}
\usepackage{graphicx} %Bilder einfügen
%\usepackage{pdfpages} %PDF einfügen
\usepackage[a4paper, margin=1in]{geometry}
\usepackage[right]{eurosym} %Euro-Zeichen
\usepackage{amssymb}
\usepackage{cite} %Quellenangaben
\usepackage{setspace} % Zeilenabstand
\usepackage[ 
   colorlinks,        % Links ohne Umrandungen in zu wählender Farbe 
   linkcolor=black,   % Farbe interner Verweise 
   filecolor=black,   % Farbe externer Verweise 
   citecolor=black,   % Farbe von Zitaten 
   urlcolor=blue	  % Farbe von Links
   ]{hyperref} %Verlinkungen
\usepackage{bookmark}
\usepackage[figure]{hypcap}
\usepackage[ngerman]{translator}
\usepackage{blindtext} % Lorem-Ipsum-Plugin
\usepackage{scrhack}
%\usepackage[
	%nonumberlist, %keine Seitenzahlen anzeigen
	%acronym,      %ein Abkürzungsverzeichnis erstellen
	%toc,          %Einträge im Inhaltsverzeichnis
	%section      %im Inhaltsverzeichnis auf section-Ebene erscheinen
	%]
%{glossaries}
\usepackage{outlines}

\usepackage{listings,xcolor} %Codeanzeige
\usepackage[normalem]{ulem}
\useunder{\uline}{\ul}{}


\usepackage{mathtools}
\DeclarePairedDelimiter\abs{\lvert}{\rvert}%
\DeclarePairedDelimiter\norm{\lVert}{\rVert}%
% Swap the definition of \abs* and \norm*, so that \abs
% and \norm resizes the size of the brackets, and the 
% starred version does not.
\makeatletter
\let\oldabs\abs
\def\abs{\@ifstar{\oldabs}{\oldabs*}}
%
\let\oldnorm\norm
\def\norm{\@ifstar{\oldnorm}{\oldnorm*}}
\makeatother

\usepackage{chngcntr}
\counterwithout{figure}{chapter}
\counterwithout{table}{chapter}



%%%%%%%%%%%%%%%%%%%%%%%%%
%%
%% Code Definition
%%
\lstdefinelanguage{powershell}%
  {
	alsodigit = {-},   
   morekeywords={abstract,break,case,catch,const,continue,do,else,elseif,%
      end,export,false,for,function,immutable,import,importall,if,in,%
      macro,module,otherwise,quote,return,switch,true,try,type,typealias,%
      using,while,new-object,psobject,Get-Command,Process,Get-ADUser,Select,ActiveXObject},%
   sensitive=true,%
   alsoother={\$},%
   morecomment=[l]\#,%
   morecomment=[n]{\#=}{=\#},%
   morestring=[s]{"}{"},%
   morestring=[m]{'}{'},%
}[keywords,comments,strings]%

%\lstset{
%	language         = [Sharp]C,
%	basicstyle       = \ttfamily,
%    numbers=left, 
%    numberstyle=\tiny, 
%    numbersep=5pt,
%    breaklines=true,
%    frame=single,
%    escapeinside={(*@}{@*)}, %nicht anzuzeigende Ausdrücke, z.B. für Labels
%    language=sh,
%    basicstyle=\ttfamily\fontsize{10}{12}\selectfont,
%    keywordstyle    = \color{dkblue},
%    stringstyle     = \color{red},
%    identifierstyle = \color{black},
%    commentstyle    = \color{gray},
%    emph            =[1]{php},
%    emphstyle       =[1]\color{black},
%    emph            =[2]{if,and,or,else,Parameter,-Property},
%    emphstyle       =[2]\color{dkyellow},
%    } 

\usepackage{color}
\definecolor{bluekeywords}{rgb}{0,0,1}
\definecolor{greencomments}{rgb}{0,0.5,0}
\definecolor{redstrings}{rgb}{0.64,0.08,0.08}
\definecolor{xmlcomments}{rgb}{0.5,0.5,0.5}
\definecolor{types}{rgb}{0.17,0.57,0.68}

\usepackage{listings}

\lstset{language=csh,
captionpos=b,
%numbers=left, %Nummerierung
%numberstyle=\tiny, % kleine Zeilennummern
frame=single, % Kasten, lines alternativ
showspaces=false,
showtabs=false,
breaklines=true,
showstringspaces=false,
breakatwhitespace=true,
escapeinside={(*@}{@*)},
commentstyle=\color{greencomments},
morekeywords={partial, var, value, get, set},
keywordstyle=\color{bluekeywords},
stringstyle=\color{redstrings},
basicstyle=\ttfamily\small,
captionpos=b,
numbers=left, 
numberstyle=\tiny, 
numbersep=5pt
}

%%%%%%%%%%%%%%%%%%%%%%%%%%%%%%%%%%%%%%%%%%%%%%%%%%%%%
%%%%%%%%%%% Sonderformatierung
%%%%%%%%%%%%%%%%%%%%%%%%%%%%%%%%%%%%%%%%%%%%%%%%%%%%%

% Seitenabstände definieren
\geometry{verbose,tmargin=3cm,bmargin=2cm,lmargin=3cm,rmargin=3cm} 

% Hurenkinder und Schusterjungen verhindern (Ja, das heißt wirklich so!!!)
\clubpenalty = 10000 \widowpenalty = 10000 \displaywidowpenalty = 10000 

\newcommand{\footfigref}[1]{\footnote{Abb. \ref{#1} auf Seite \pageref{#1}}}

%% Bei Referenzen im Text wird jetzt bei allen Ebenen "Kapitel" vorgestellt, z.b. Kapitel 2, Kapitel 2.2, Kapitel 6.3.2
\addto\extrasngerman{%
    \def\sectionautorefname{Kapitel}%
    \def\subsectionautorefname{Kapitel}%
    \def\subsubsectionautorefname{Kapitel}%
    }

% Vertikaler Abstand zwischen Ende Textblock - Ende Fußzeile --> Abstand der Seitenzahl von Rand erhöhen 
\setlength{\footskip}{10mm}

% Abstand vor/nach Überschriften verändern

\RedeclareSectionCommand[%
    beforeskip=0.5\baselineskip,
    afterskip=0.5\baselineskip
]{chapter}

\RedeclareSectionCommand[%
    beforeskip=0.5\baselineskip,
    afterskip=0.5\baselineskip
]{section}

\RedeclareSectionCommand[%
    beforeskip=0.1\baselineskip,
    afterskip=0.1\baselineskip
]{subsection}

\RedeclareSectionCommand[%
    beforeskip=0.01\baselineskip,
    %%afterskip=0.2\baselineskip
]{paragraph}

\usepackage[acronym, nonumberlist]{glossaries} %% use after hyperref %Glossar-Paket laden


\setlength{\abovecaptionskip}{4pt}  % 1pc=12pt 
\setlength{\belowcaptionskip}{0pt}
%\setlength{\textfloatsep}{4pt}
\setlength{\intextsep}{1pc}

%% Verkleinerung der Textgröße unter Abbildungen
\addtokomafont{caption}{\small}

% falsche Default-Silbentrennung überschreiben
%falsche Default-Silbentrennung überschreiben
\hyphenation{Soft-ware-ent-wick-lung}

% Den Punkt am Ende der Glossareinträge deaktivieren
\renewcommand*{\glspostdescription}{}

%Glossar-Befehle anschalten
%\makeglossaries
\makenoidxglossaries

% sorgt dafür, dass bei Leerzeile die Einrückung verhindert und stattdessen eine Leerzeile eingefügt wird % erspart bigskips und erhöht die Lesbarkeit im LaTeX-Text 
\KOMAoptions{parskip=full*}

% ändert Titelschriftart in Serifen-Normalschriftart
\addtokomafont{disposition}{\rmfamily} 



\loadglsentries{glossar.tex}

%%%%%%%%%%%%%%%%%%%%%%%%%%%%%%%%%%%%%%%%%%%%%%%%%%%%%
%%%%%%%%%%% Textbausteine
%%%%%%%%%%%%%%%%%%%%%%%%%%%%%%%%%%%%%%%%%%%%%%%%%%%%%
%%%%%%%%%%%% Studentennamen
\newcommand{\studentNameEins}{Jahn, Marko}
\newcommand{\studentNameZwei}{Schenkewitz, Mario}
\newcommand{\studentNameDrei}{Zilius, Sven}
%%%%%%%%%%%% Typ der Arbeit
\newcommand{\type}{Studienprojekt I}
%%%%%%%%%%%% Thema
\newcommand{\topic}{Schatzinsel}
%%%%%%%%%%%% Untertitel
\newcommand{\subtopic}{Entwickeln eines Spiels durch objektorientiertes Programmieren}
%%%%%%%%%%%% Studienbereich
\newcommand{\studienbereich}{Technik}
%%%%%%%%%%%% Fachrichtung
\newcommand{\fachrichtung}{Duales Studium Wirtschaft • Technik}
%%%%%%%%%%%% Studiengang
\newcommand{\studiengang}{Informatik}
%%%%%%%%%%%% Betrieb
\newcommand{\company}{}
%%%%%%%%%%%% Betreuer HWR
\newcommand{\betreuerHS}{Zimmermann, Arthur}
%%%%%%%%%%%% Betreuer Unternehmen
\newcommand{\betreuerUnt}{}
%%%%%%%%%%%% Jahrgang
\newcommand{\jahrgang}{2020}
%%%%%%%%%%%% Matrikelnummer
\newcommand{\MatrEins}{}
\newcommand{\MatrZwei}{685593}
\newcommand{\MatrDrei}{}
%%%%%%%%%%%% Modul
\newcommand{\Modul}{IT3161, Studienprojekt I}
%%%%%%%%%%%% Anzahl Wörter
\newcommand{\wordsnr}{TDB}
%%%%%%%%%%%% Fertigstellungsdatum
\newcommand{\Dateofcompletion}{TBD}
%%%%%%%%%%%%%%%%%%%%%%%%%%%%%%%%%%%%%%%%%%%%%%%%%%%%%>>>>>>>


\begin{document}

%%%%%%%%%%%%%%%%%%%%%%%%%%%%%%%%%%%%%%%%%%%%%%%%%%%%%>>>>>>>
%%%%%%%%%%% Titelblatt

%% Anordnung und Aussehen von Titel und Untertitel

\subject{\type}

\title{
\normalfont\endgraf\rule{\textwidth}{.4pt}
\begingroup
	\centering
	\linespread{1.5}
	\huge\topic
\endgroup
\linespread{1}
\ \\ % Falls kein Subtopic, auskommentieren
\ \\ % Falls kein Subtopic, auskommentieren
\large\subtopic % Falls kein Subtopic, auskommentieren
\endgraf\rule{\textwidth}{.4pt}
}
 
%%Eigentlich nicht besonders schön, aber Koma erlaubt die Anordnung eines weiteren Felden (hier: Fachbereich) nicht
\date{\normalsize fertiggestellt am \Dateofcompletion \\ \textbullet \\ Fachbereich Duales Studium Wirtschaft / Technik \\
Hochschule für Wirtschaft und Recht Berlin}
%% \date muss leer angegeben werden, um die Default-Datumsanzeige zu unterdrücken

\publishers{
	\begin{tabular}{l l}
	\textbf{\normalsize{}} & \normalsize{}  \tabularnewline
	%\textbf{\normalsize{}} & \normalsize{}  \tabularnewline
	\textbf{\normalsize{Name:}} & \normalsize{\studentNameEins, \MatrEins} \tabularnewline
	\textbf{\normalsize{Name:}} & \normalsize{\studentNameZwei, \MatrZwei} \tabularnewline
	\textbf{\normalsize{Name:}} & \normalsize{\studentNameDrei, \MatrDrei} \tabularnewline
	\textbf{\normalsize{Fachrichtung:}} & \normalsize{\fachrichtung} \tabularnewline
	\textbf{\normalsize{Studiengang:}} & \normalsize{\studiengang} \tabularnewline
	\textbf{\normalsize{Studienjahrgang:}} & \normalsize{\jahrgang} \tabularnewline
	\textbf{\normalsize{Betreuer HS:}} & \normalsize{\betreuerHS} \tabularnewline
	\textbf{\normalsize{Anzahl der Wörter:}} & \normalsize{\wordsnr}\\
	
%	\normalsize Vom Betreuer zur Kenntnis genommen:\vspace{2em}\\
%	\end{tabular}
%	\normalsize
%	\begin{tabular}{lp{10em}l} 
% 	\hspace{1cm}   && \hspace{4cm} \\\cline{1-1}\cline{3-3} 
% 	Ort, Datum     && \betreuerHS
	\end{tabular}
	}

\titlehead{\begin{center}
    \includegraphics[scale=1]{bilder/HWR-Logo.jpg}
    \end{center}
    }

\maketitle

\onehalfspacing % anderthalbfacher Zeilenabstand


%%%%%%%%%%%%%%%%%%%%%%%%%%%%%%%%%%%%%%%%%%%%%%%%%%%%%%%%%%%%%%%%%%%%%%%%%%%%%%%%%%%%%%%%%%%%%%%%%%%%%%%%%%%%%%%%%%%%%%%%%%%
%%%%%%%%%%% Dokumenteninhalt START
%%%%%%%%%%%%%%%%%%%%%%%%%%%%%%%%%%%%%%%%%%%%%%%%%%%%%%%%%%%%%%%%%%%%%%%%%%%%%%%%%%%%%%%%%%%%%%%%%%%%%%%%%%%%%%%%%%%%%%%%%%%

%%%%%%%%%%%%%%%%%%%%%%%%%%%%%%%%%%%%%%%%%%%%%%%%%%%%%
%%%%%%%%%%% Abstract
\chapter*{Abstract}
\addcontentsline{toc}{chapter}{Abstract}


\pagenumbering{Roman} % römische Seitenzahlen

\chapter*{Kurzfassung}
\addcontentsline{toc}{chapter}{Kurzfassung}


%%%%%%%%%%%%%%%%%%%%%%%%%%%%%%%%%%%%%%%%%%%%%%%%%%%%%
%%%%%%%%%%% Inhaltsverzeichnis, Tabellen, Abbildungen, etc.
\newpage

\tableofcontents{}
\addcontentsline{toc}{chapter}{Inhaltsverzeichnis}

\listoffigures
%\listoftables

%\section*{Hinweis}
%
%Aus Gründen der besseren Lesbarkeit wird im Text verallgemeinernd die weibliche Form verwendet. Diese Formulierungen umfassen gleichermaßen weibliche und männliche Personen.
%\clearpage

\addcontentsline{toc}{chapter}{Akronyme}
\printnoidxglossaries

\clearpage

%% arabische Seitenzahlen
\pagenumbering{arabic}

%%%%%%%%%%%%%%%%%%%%%%%%%%%%%%%%%%%%%%%%%%%%%%%%%%%%%
%%%%%%%%%%% Einführung

\chapter{Einleitung}\label{sec:einleitung}
Die \gls{OOP} ist ein heutzutage weit verbreitetes Computerprogrammiermodell oder Programmierparadigma. In der \gls{OOP} werden Anwendungen um so genannte Objekte und deren Daten herum strukturiert und nicht um Funktionen oder Logik. So fokussiert sich die \gls{OOP} auf Objekte mit denen eine Anwendung interagiert, diese Objekte können als Datenfeld mit Attributen und eigenem Verhalten angesehen werden.

Einen tieferen Einblick in das Thema der \gls{OOP} soll der Abschnitt \ref{sec:Gundlagen} bieten, in dem die Basis für die \gls{OOP} näher erörtert wird.

Thema dieses Studienprojektes ist es sich dieses Programmierparadigma zunutze zu machen, um ein Spiel zu entwickeln.\\
Der Titel des Spiels ist \glqq Schatzinsel\grqq, es handelt von drei Piraten die sich auf einer Insel befinden, auf der ebenfalls ein Schatz versteckt ist. Es ist ihnen unbekannt wo sich dieser Schatz befinden. Durch zufällige Bewegung über die Insel wollen sie versuchen diesen Schatz zu finden. Jeder der Piraten hat verschiedene Eigenschaften, wie zum Beispiel Geschwindigkeit.\\
Sobald einer der Piraten den Schatz entdeckt, ruft er die anderen zu sich. Am Ende versammeln sie sich alle um den Schatz.

Dieser Umstand soll durch das Spiel modelliert und gegebenenfalls erweitert werden.

\chapter{Grundlagen}\label{sec:Gundlagen}
\section{Objektorientierte Programmierung}
\gls{OOPL} eignen sich durch die Konzentration auf die Objekte mit denen das Programm interagieren soll dazu, dass Software Wartbar und lesbar bleibt. Objekte und ihre Methoden können getrennt voneinander erstellt und programmiert werden. Diese Unterteilung in Einzelteile ermöglicht des Weiteren die Wiederverwendbarkeit von Code, also dem Quelltext, entweder in dem die Klasse selbst wiederverwendet wird, oder eine andere Klasse von ihr erbt. Dazu mehr im Abschnitt \ref{sec:Vererbung}

\gls{OOP} ist ein Programmierstil, der mittlerweile in Unternehmen, Industrie und akademischen Kreisen seinen Platz gefunden hat. Es wird ein natürlicheres und dadurch ein mächtigeres Entwurfsparadigma insofern erwartet, dass ein Programmierer produktiver ist, wenn das Anwendungsmodell näher am Problem modelliert werden kann, dass gelöst wird.\\
Anders gesagt ist eine \gls{OOPL} mächtiger, weil es einem Programmierer erlaubt ein Problem zu lösen, dass die Struktur einer Welt nachbildet in welcher das Problem entstand. Vertrautheit mit solch einer Sprache eröffnet Wege zu einer klaren und natürlichen Lösung \cite{OOPL}.

Es werden also im betrachteten Fall der Schatzinsel die Objekte, die aus der echten Welt bekannt sind, so modelliert, dass die Interaktionen und Verhaltensweisen natürlich in einer Anwendung modelliert werden können.

\section{Datenabstraktion und Verkapselung}
Einer der ersten Schritte bei der \gls{OOP} ist es alle Objekte zu sammeln, die manipuliert werden müssen. Im Fall des Spiels \glqq Schatzinsel\grqq \hspace{0.1cm}gibt es zum Beispiel Piraten, eine Insel, ein Schatz, aber auch eine graphische Benutzeroberfläche die alle als Objekt betrachtet werden können. Sie müssen untereinander aber auch mit dem Benutzer interagieren können.

Einer der großen Durchbrüche bei der Gestaltung von Programmiersprachen geschah, als es ermöglicht wurde, dass Programmierer eigene Datentypen definieren konnten. C.A.R. Hoare's record class in Simula67 gilt als die Quelle der Idee.\cite{OOPL}
Verbreitung des, durch Programmierer definierten, Datentyps fand durch die Sprache Pascal statt. \cite{Pascal}

So kann in Pascal ein Verbund definiert werden, der selbst aus Daten besteht, siehe Listing \ref{lst:record1}. 

\begin{lstlisting}[language=Pascal, caption=Verbund Beispiel Definition \cite{PascalVerbund}, label={lst:record1}]
TYPE Auto = RECORD
KFZ: String[12];
Gewicht: Real;
AnzPassagiere: Integer;
END;
\end{lstlisting}

Dieses Beispiel (Listing \ref{lst:record2}) definiert einen neuen Datentyp \textit{Auto}, in dem jeweils die Art des KFZ, das Gewischt und die Anzahl der Passagiere angegeben wird. Nun kann man den Typen \textit{Auto} nutzen, um beliebig viele Daten-Entitäten zu erstellen.

\begin{lstlisting}[language=Pascal, caption=Verbund Beispiel Instantiierung \cite{PascalVerbund}, label={lst:record2}]
VAR a : Auto;
BEGIN
a.KFZ := 'B XY 123456';
a.Gewicht := 555.77;
a.AnzPassagiere := 5;
END;
\end{lstlisting}

Manipulationen der Daten innerhalb des Verbunds werden extern vom Datentyp selbst betrachtet, die Syntax der Sprache indiziert somit nur eine Gruppierung von Daten.

Verbunde in Pascal werden homogen genannt, wenn alle Daten innerhalb des Verbunds denselben Typ haben. Sie sind heterogen, wenn es verschiedene Typen, wie im Beispiel, gibt.

Pascal erlaubt zwar die Erstellung von heterogenen Verbunden, jedoch ist es nicht möglich in diesen Verbunden Operationen einzubauen. Anders als bei Pascal ist die bei Java zum Beispiel möglich.  

\begin{lstlisting}[language=Java, caption=Operation innerhalb einer Klasse \cite{OOPL}, label={lst:Java-Operation}]
class Point {
private int x, y;
	public int getX() { return x; }
	public int getY() { return y; }
	public void setXY(int x, int y) {
		if (x >= 0 && y >= 0) {
			this.x = x;
			this.y = y;
		}
	}
}
\end{lstlisting}

Im Listing \ref{lst:Java-Operation} wird zum Beispiel durch die Funktion \glqq setXY()\grqq{}  ermöglicht, dass x und y gesetzt werden können. Es ist zu bemerken, dass verschiedene Sichtbarkeiten gesetzt werden. Dadurch, dass x und y auf \emph{private} gesetzt sind, können sie also von außen nur durch die Funktion \glqq setXY()\grqq{} geändert werden.  
Dies veranschaulicht auch das Konzept der Verkapselung \cite{OOPL} Seite 4.  

Datenabstraktion in ihrer vollen Kapazität beschreibt also die Möglichkeit für Programmierer ihre eigenen Datentypen zu definieren, diese Typen können eine Heterogene Struktur haben und die Sicht auf diese kann eingeschränkt werden. Außerdem ist es möglich Operationen innerhalb dieser Datentypen zu haben.  

In diesem Zusammenhang kann gesagt werden, dass Objekte also aus Attributen und Operationen, auch Funktionen oder Methoden genannt, bestehen und Instanzen von diesen Datentypen sind. Diese Datentypen werden im Allgemeinen auch als Klassen bezeichnet.


\section{Vererbung}\label{sec:Vererbung}
Die Beziehung zwischen Typen in spielt eine wichtige Rolle in den \gls{OOPL}, so ist die Vergleichbarkeit von Typen wichtig, um Objekte zu vergleichen.  
Diese Beziehungen werden durch Vererbung erweitert und bereichert. Es stellt sich ein System aus Vererbungs- und Subtypenbeziehungen dar, diese Tragen zur Komplexität einer Sprache bei. Dabei spielen verschiedene Begriffe zentrale Rollen.  
Zum Beispiel wird gesagt, dass Klasse B von Klasse A \emph{erbt} oder \emph{abgeleitet} ist, wenn Objekte der Klasse B mindestens dieselben Attribute und Verhaltensweisen wie Objekte der Klasse A haben. Die Objekte der Klasse B können auch mit eigenen Attributen und Methoden erweitert werden.  
B \emph{erbt} diese Attribute von A. Dann ist A die \emph{Basisklasse}, manchmal auch als \emph{Elternklasse} bezeichnet, von der B abgeleitet wird.  
Im englischen werden oft auch die Begriffe \emph{Superklass} (Oberklasse) und \emph{Subclass} (Unterklasse) verwendet.  
Subtypenbeziehungen gestalten sich ähnlich zur Vererbung. Damit eine Subtypenbeziehung besteht, müssen Objekte des Typs B auch als Objekte des Typs A gelten. Nicht nur die Attribute des Objekts vom Typ B stimmen mit denen vom Typ A überein, auch die Zugriffsrechte und Sichtbarkeit auf die verschiedenen Attribute müssen übereinstimmen. Das heißt, dass alle Attribute und Methoden, die in Typ A zur Verfügung stehen, auch in Typ B zur Verfügung stehen.  
Wenn diese Bedingung erfüllt ist, dann können auch immer Objekte vom Typ B als Argumente an Parameter des Typen A gegeben werden. Im englischen wird gesagt B ist ein \emph{subtype} (Untertyp) von A beziehungsweise A ist ein \emph{supertype} (Obertyp) von Typ B. Im Allgemeinen ist es so, dass es immer, wenn die Funktionalität der Vererbung gegeben ist, auch immer die Funktionalität für Subtypen in Programmiersprachen existiert.  
  
Ein Datentyp, der Vererbung eingearbeitet hat, wird auch als Klasse bezeichnet (s.O.).

\begin{lstlisting}[language=Java, caption=Erweiterung der Point Klasse aus Listing \ref{lst:Java-Operation} \cite{OOPL}, label={lst:Vererbung}]
class Disk extends Point {
	private double radius;
	private Color color;
	public double getRadius() { return radius; }
	public void setRadius(double radius) {
		if (radius > 0) this.radius = radius;
		}
	public Color getColor() { return color; }
	public void setColor(Color color) { this.color = color; }
}
\end{lstlisting}

Im Listing \ref{lst:Vererbung} wird die Klasse \glqq Point\grqq{}  aus \ref{lst:Java-Operation} erweitert. Da eine Scheibe einen Punkt als Basis hat, kann die Klasse auch so erweitert werden, dass man einen Radius und z.B. eine Farbe hinzufügt. Die Vererbung selber ändert nichts an der Verwendung dieser neuen Klasse. Sie muss immer noch Instantiiert werden und über Methoden ihre Parameter gesetzt werden, dafür dienen die Methoden setRadius(), SetColor() und die geerbte Methode setXY(). Es ist also Beispielweise möglich mit den folgenden Befehlen ein entsprechendes Objekt zu erstellen siehe Listing \ref{lst:Instantiierung}.

\begin{lstlisting}[language=Java, caption=Instantiierung der Klasse Disk \cite{OOPL}, label={lst:Instantiierung}]
Disk d = new Disk();
d.setXY(100,210);
d.setRadius(25.7);
d.setColor(Color.red);
\end{lstlisting}

Es ist zu bemerken, dass die Methoden über das Objekt aufgerufen werden. Es gibt die Möglichkeit auch Methoden auszuführen, die über Klassen ausgeführt werden. Solche Methoden nennt man auch statische Methoden, sie werden über das Schlüsselwort \emph{static} im Quelltext der Sprachen Java, C++ und C\# identifiziert.

\section{Polymorphismus}
%Seite 32
Subklassen sind oft spezialisierte klassen
daher notwendig , dass ein spezialisiertes Objekt ihre Methoden anders ausführt
wir benötigen die Möglichkeit Methoden zu überschreiben \emph{override}

Subklassen sind oft spezialisierte Klassen, die von einer Oberklasse erben. Oft ist es zu erwarten, dass eine solche spezialisierte Klasse sich anders verhält, wenn man eine Methode der Oberklasse aufruft. Man kann zur Veranschaulichung eine Klasse betrachten, die ein beliebiges Tier darstellen soll siehe Listing \ref{lst:Tier}.

\begin{lstlisting}[language=Java, caption=Klasse Tier, label={lst:Tier}]
class Tier
{
    private String name;
    private short alter;
 
    public void gibLaut()
    {
        System.out.println("Geraeusche eines allgemeinen Tiers");
    }

}
\end{lstlisting}

Wenn man nun Unterklassen von der Klasse Tier ableitet, könnte man erwarten, dass die Methode \emph{gibLaut()}, je nachdem, um was für ein Tier es sich handelt, eine andere Ausgabe verursacht.  
Es ist also notwendig eine Möglichkeit zu haben die Verhaltensweisen und Methoden von Oberklassen zu überschreiben. Diese veränderbarkeit der Implementation von Funktionen nennt man auch Polymorphismus. 
So kann man in der Programmiersprache Java mit der Deklaration \emph{\@Override} eine Funktion überschreiben siehe \ref{lst:Katze}.

\begin{lstlisting}[language=Java, caption=Klasse Katze, label={lst:Katze}]
class Katze extends Tier
{
    @Override
    public void gibLaut()
    {
        System.out.println("Miau");
    }
}
\end{lstlisting}

Je nach Sprache ist diese Funktionalität anders implementiert. In der Sprache C++ werden Funktionen, die überschreibbar sein sollen mit dem Schlüsselwort \emph{virtual} gekennzeichnet siehe \ref{lst:TierC++}.

\begin{lstlisting}[language=C++, caption=Tier und Katzen Beispiel in C++, label={lst:TierC++}]
class Tier {
	public:
	virtual char * gibLaut() { return "Geraeusche eines allgemeinen Tiers"; }
	};
	
class Katze : public Tier {
	public:
	char * gibLaut() { return "Miau"; }
};
\end{lstlisting}

\section{Überladen}
Überladen bezieht sich auf Funktionen und Operatoren und bedeutet, dass diese, unter der Verwendung des gleichen Namens, in Abhängigkeit der Typen der übergebenen Parameter verschieden ausgeführt werden. Betrachtet man zum Beispiel die Addition als Operator, so wird bei der Berechnung der Summe von zwei Zahlen je nach Typ der Zahlen die Operation anders ausgeführt.  
Wenn man Zahlen von verschiedenen Typen addieren möchte, dann wird oft impliziert eine Umwandlung der Typen vollzogen. Zum Beispiel wenn eine ganze Zahl mit einer Fließkommazahl addiert wird, dann wird die ganze Zahl in eine Fließkommazahl umgewandelt, bevor sie addiert werden.  
Der + Operator für die Addition wird in einigen Sprachen auch so überladen, dass mit ihm Text verkettet werden kann.  
  
Es ist Programmierern somit auch möglich selber Operatoren oder Funktionen zu überladen.

\chapter{Umgebungen}\label{sec:Umgebungen}
\section{Java Swing}
Java Swing ist eine \gls{GUI} Bibliothek für Java. Es gehört zu den Java Foundation Classes, eine Sammlung von Bibliotheken für dir Programmierung von graphischen Benutzerschnittstellen. Andere Bibliotheken in dieser Sammlung sind zum Beispiel Java 2d, Java Accessibility API, Drag-and-Drop-API und das Abstract Window Toolkit (AWT) \cite{JFC}. Swing baut auf AWT auf und ist mit den anderen APIs verwoben.  
Es gibt auch Konkurrenten zu Swing, z.B. das für Eclipse entwickelte SWT. Auf dieses wird aber nicht eingegangen, da Swing auch mit anderen Entwicklerwerkzeugen und -umgebungen Kompatibel ist.  
Swing ist, im Gegensatz zu AWT, Plattformunabhängig. Außerdem ist es modular und objektorientiert aufgebaut. Die Plattformunabhängigkeit wird dadurch gewährleistet, dass alle Swing-Komponenten direkt von Java gerendert werden. Nachteil davon ist, dass Swing Anwendungen nicht wie Anwendungen aussehen, die nativ auf dem darunterliegenden Betriebssystem entwickelt wurden.  
Es gibt so genannte \glqq Look-and-Feels\grqq{}, die das Aussehen und Verhalten modifizieren.

\section{Unreal Engine}
\section{Unity Engine}


\chapter{Verwendete Werkzeuge}
\section{Git und Github}
\section{Eclipse}
\section{Unreal Engine Editor}
\section{Unity Engine Editor}



\chapter{Realisierung}\label{sec:Realisierung}
\section{Zusammenfassung der Anforderungen}
Wie in \ref{sec:einleitung} Beschrieben sind für dieses Projekt einige Anforderungen gegeben, die Umgesetzt werden sollen. Für die Betrachtung der Entwicklung dieses Projekts in den verschiedenen Umgebungen werden also folgende Anforderungen festgelegt.
Diese Funktionen werden nach den \glqq MASTeR - Die SOPHIST-Templates für Requirements\grqq{} von Sophist formuliert\cite{sophist}.

\textbf{Anforderungen}
\textbf{Spielstart}
\begin{itemize}\vspace{-1em}
\setlength{\itemsep}{-1em}
	\item Wenn das Programm gestartet wird, muss das Spiel 1 Spielfeld, 1 Schatz und 3 Piraten Objekte erstellen.
	\item Wenn die Piraten erstellt werden, muss das Spiel ihnen verschiedene Geschwindigkeiten zuordnen.
	\item Wenn das Programm gestartet wird, muss das Spiel die Graphische Oberfläche darstellen.
	\item Wenn das Programm aktualisiert muss das Spiel die Piraten \textbf{zufällig} über das Spielfeld bewegen.
\end{itemize}

\textbf{Während des Spiels}
\begin{itemize}\vspace{-1em}
\setlength{\itemsep}{-1em}
	\item Wenn sich ein Pirat bewegt, sollte das Spiel ein Feld wählen, dass er in den letzten 5 Positionen nicht betreten hat.
	\item Wenn das Spiel gestartet ist, muss das Spiel eine Option für die Sichtbarkeit des Schatzes anbieten.
	\item Wenn ein Pirat in der Nähe des Schatzes ist, muss das Spiel den Schatz sichtbar machen, wenn er unsichtbar ist.
	\item Wenn ein Pirat in der Nähe des Schatzes ist, muss das Spiel ihn Richtung Schatz bewegen. (Nicht mehr zufällig.)
	\item Wenn ein Pirat in der Nähe des Schatzes ist, muss das Spiel die anderen Piraten in Richtung des Schatzes bewegen.
\end{itemize}

\section{Umsetzung in Java}

\section{Umsetzung in Unreal Engine}

\section{Umsetzung in Unity Engine}

%\section{Anforderung 1}
%\subsection{Ansatz}
%\subsection{Umsetzung}
%
%\section{Anforderung 2}
%\subsection{Ansatz}
%\subsection{Umsetzung}
%
%\section{Anforderung 3}
%\subsection{Ansatz}
%\subsection{Umsetzung}
%
%\section{Anforderung 4}
%\subsection{Ansatz}
%\subsection{Umsetzung}

\chapter{Ergebnisse}

\chapter{Fazit}


\chapter{Ausblick}



%%% Codesnippets

%% Liste
%\textbf{Grundwerte}
%\begin{itemize}\vspace{-1em}
%\setlength{\itemsep}{-1em}
%	\item \textbf{Integrität}: Die Daten oder das System sind nicht von Unbefugten veränderbar. Alle Änderungen müssen nachvollziehbar sein.
%	\item \textbf{Vertraulichkeit}: Die Daten dürfen nur von autorisierten Personen abgerufen und verändert werden. Dies gilt sowohl für gespeicherte Dateien, als auch die Datenübertragung.
%	\item \textbf{Verfügbarkeit}: Die Daten und Systeme müssen erreichbar sein. Systemausfälle müssen vermieden werden. Wenn man eine Verfügbarkeit von "four nines" (99.99\% der Zeit) gewährleisten möchte, dann sind z.B. nur Ausfälle für 52:36 Minuten pro Jahr vertretbar.
%\end{itemize}

%% Bild
%\begin{figure}[H]
%  \centering
%  \includegraphics[scale=0.5]{bilder/Netzplan.JPG}
%  \label{fig:kernabsicherung}       %fig:ID
%  \caption[Beispielhafter Netzplan nach BSI 200-2]{\cite{BSI200_2} Seite 90: Auszug aus dem bereinigten Netzplan der RECPLAST GmbH (Teilausschnitt)}    %Bildunterschrift
%\end{figure}


%% Code
%\begin{lstlisting}[language=csh, caption=Python example]
%public async Task<IActionResult> Archiv()
%    {
%        string UserPersId = User.Identity?.Name ?? "Empty";
%
%        return View(await _context.Dokumente.Where(x => x.PersonalNummer == UserPersId).ToListAsync());
%    }
%//test
%\end{lstlisting}


%%%%%%%%%%%%%%%%%%%%%%%%%%
% Quellen
%%%%%%%%%%%%%%%%%%%%%%%%%

\bibliography{literatur}

\bibliographystyle{hwrbib}
%% \bibliographystyle{alpha} %% tu es nicht, niemals, das ist eklig, nicht einkommentieren

\chapter*{Ehrenwörtliche Erklärung}
\addcontentsline{toc}{chapter}{Ehrenwörtliche Erklärung}

% Keine Kopf- und Fußzeilen ausgeben
\thispagestyle{empty}
% Aber trotzdem ins Inhaltsverzeichnis aufnehmen
%\addcontentsline{toc}{section}{Eidesstattliche Erklärung}

% Hier der offizielle Text der eidesstattlichen Erklärung
Ich erkläre ehrenwörtlich:
\begin{enumerate}
	\item dass ich meine Bachelor-Thesis selbstständig verfasst habe,
	\item dass ich die Übernahme wörtlicher Zitate aus der Literatur sowie die Verwendung der Gedanken anderer Autoren an den entsprechenden Stellen innerhalb der Arbeit gekennzeichnet habe,
	\item dass ich meine Bachelor-Thesis bei keiner anderen Prüfung vorgelegt habe.
\end{enumerate}
Ich bin mir bewusst, dass eine falsche Erklärung rechtliche Folgen haben wird.
% Etwas Abstand für die Unterschrift
\vspace{2cm}

% Hier kommt die Unterschrift drüber
\begin{tabular}{lp{4em}l} 
 \hspace{5cm}   && \hspace{6cm} \\\cline{1-1}\cline{3-3} 
 Ort, Datum     && \studentNameEins \\ [10ex]
 \hspace{5cm}   && \hspace{6cm} \\\cline{1-1}\cline{3-3} 
 Ort, Datum     && \studentNameZwei \\ [10ex]
 \hspace{5cm}   && \hspace{6cm} \\\cline{1-1}\cline{3-3} 
 Ort, Datum     && \studentNameDrei \\ [10ex]
\end{tabular}

\end{document}

